\documentclass{article}

% Recommended, but optional, packages for figures and better typesetting:
\usepackage{microtype}
\usepackage{graphicx}
\usepackage{subcaption}
\usepackage{booktabs} % for professional tables
\usepackage{hyperref}
\newcommand{\theHalgorithm}{\arabic{algorithm}}
\usepackage{icml2026}

\usepackage{amsmath}
\usepackage{amssymb}
\usepackage{mathtools}
\usepackage{amsthm}
\usepackage[capitalize,noabbrev]{cleveref}

%%%%%%%%%%%%%%%%%%%%%%%%%%%%%%%%
% THEOREMS
%%%%%%%%%%%%%%%%%%%%%%%%%%%%%%%%
\theoremstyle{plain}
\newtheorem{theorem}{Theorem}[section]
\newtheorem{proposition}[theorem]{Proposition}
\newtheorem{lemma}[theorem]{Lemma}
\newtheorem{corollary}[theorem]{Corollary}
\theoremstyle{definition}
\newtheorem{definition}[theorem]{Definition}
\newtheorem{assumption}[theorem]{Assumption}
\theoremstyle{remark}
\newtheorem{remark}[theorem]{Remark}

\icmltitlerunning{Wasserstein Spectral Alignment}

\begin{document}

\twocolumn[
  \icmltitle{Wasserstein Spectral Alignment: A Principled Optimal Transport Framework for Model Merging}

  \begin{icmlauthorlist}
    \icmlauthor{Theorist Agent}{equal,yyy}
  \end{icmlauthorlist}

  \icmlaffiliation{yyy}{Department of Computer Science, Autonomous Research Institute}
  \icmlcorrespondingauthor{Theorist Agent}{theorist@research.org}

  \icmlkeywords{Model Merging, Optimal Transport, Bures-Wasserstein, Spectral Collapse}

  \vskip 0.3in
]

\printAffiliationsAndNotice{}

\begin{abstract}
  Model merging offers an efficient alternative to multi-task learning by combining independently fine-tuned models without additional training. However, prevailing methods like Task Arithmetic suffer from severe spectral distortion and destructive interference of parameter updates, leading to a ``spectral collapse'' that degrades downstream performance. While empirical fixes like isotropic scaling have been proposed, they lack a unified mathematical foundation to align and preserve the underlying statistical structure of fine-tuned weights. In this paper, we introduce \textbf{Wasserstein Spectral Alignment (WSA)}, a mathematically rigorous framework that models task-specific parameter updates as multivariate probability distributions and reformulates model merging as a Bures-Wasserstein barycenter problem over their covariances. WSA aligns and re-scales the merged updates to preserve both the spectral energy and the task-specific directional correlations in a theoretically guaranteed manner. Under orthogonal constraints, we prove that WSA generalizes previous orthogonal model merging techniques. Extensive experiments across vision benchmarks (CIFAR-10, SVHN, FashionMNIST) demonstrate that WSA consistently outperforms existing baselines, achieving superior multi-task accuracy and mitigating parameter interference.
\end{abstract}

\section{Introduction}
\label{submission-introduction}

Deep neural networks are typically fine-tuned on specialized downstream datasets to excel at individual tasks. However, deploying a separate model for each task is computationally expensive. Multi-task learning (MTL) addresses this by training a single model on all tasks simultaneously, but suffers from high training costs, negative transfer, and optimization instabilities.

To address these limitations, \emph{model merging} has emerged as a promising paradigm to create a unified multi-task model by directly combining independently fine-tuned weights at the parameter level. Among these, \textbf{Task Arithmetic (TA)} \cite{ilharco22} has gained popularity due to its simplicity, alongside related techniques like Model Soups \cite{wortsman22}, Git Re-Basin \cite{ainsworth22}, and ZipIt! \cite{stoica23}. TA computes a ``task vector'' $\tau_i = W_i - W_0$ representing the difference between the fine-tuned weight $W_i$ and the pretrained weight $W_0$, and merges them via a linear combination:
\begin{equation}
  W_{\text{TA}} = W_0 + \sum_{i=1}^K \lambda_i \tau_i
\end{equation}
Despite its empirical success, linear averaging in Euclidean space ignores the complex, non-linear geometric structure of the weight manifold. When independent task updates are averaged, they often undergo destructive interference due to feature and subspace misalignment. This results in \emph{spectral collapse}, where the singular values of the merged update matrix are severely attenuated, driving the merged model back toward the pretrained base model and sacrificing downstream adaptations.

To counter this, recent works have proposed heuristic fixes. For example, \textbf{Sharpness-Aware Isotropic Merging (SAIM)} \cite{saim} uses isotropic scaling to uniformize the singular value spectrum of the merged updates. However, uniform scaling ignores the rich correlation structures and directional intensities of the task-specific weights. On the other hand, \textbf{Orthogonal Model Merging (OrthoMerge)} \cite{orthomerge} restricts parameter updates to the orthogonal group, but this constraint is highly restrictive and unsuitable for standard additive fine-tuning techniques like full-weight adaptation. Other approaches such as Fisher-Weighted Averaging \cite{matena22} and Relative Covariance Alignment (RECO) \cite{reco24} try to incorporate curvature information but can be computationally prohibitive for large-scale models.

In this work, we present a mathematically rigorous, hyperparameter-free solution to this geometric dilemma: \textbf{Wasserstein Spectral Alignment (WSA)}. We approach model merging through the lens of Optimal Transport (OT) \cite{kantorovich42, villani03, villani08, peyre19} and Riemannian geometry. By modeling the row (or column) vectors of each task-specific update matrix as empirical samples from a zero-mean multivariate Gaussian distribution, we can represent each task's parameter update by its covariance matrix $\Sigma_i$. 

We then formulate the merging of these covariances as finding their **Bures-Wasserstein Barycenter** $\Sigma^*$ \cite{agueh11}, which minimizes the sum of 2-Wasserstein distances. Unlike Euclidean averaging, the Wasserstein barycenter preserves the average spectral energy and directional correlations of the underlying distributions. We then compute the standard Task Arithmetic update $T_{\text{TA}}$ and perform a minimal spectral projection to align its covariance with the optimal barycenter $\Sigma^*$. This yields the WSA update $T_{\text{WSA}}$, which we reshape back to the original tensor dimensions.

Our contributions are summarized as follows:
\begin{itemize}
  \item We introduce **Wasserstein Spectral Alignment (WSA)**, a principled model merging framework grounded in optimal transport and Bures-Wasserstein barycenters \cite{agueh11}.
  \item We prove that WSA mathematically preserves the expected spectral energy of individual task updates and generalizes previous orthogonal Procrustes formulations under strict orthogonal constraints.
  \item We provide a highly efficient, provably convergent fixed-point algorithm \cite{bhatia19} that scales to large neural networks, requiring only small-matrix SVDs or eigendecompositions.
  \item We demonstrate empirically that WSA consistently outperforms Task Arithmetic and Isotropic Merging on image classification benchmarks, maintaining robust multi-task capabilities without any parameter tuning.
\end{itemize}

\section{Related Work}
\label{submission-related-work}

\textbf{Parameter-Space Model Merging.} Linear merging methods, such as Task Arithmetic \cite{ilharco22} and Ties-Merging \cite{yadav23}, perform arithmetic averaging of fine-tuned weights. Model Soups \cite{wortsman22} extends this by averaging multiple fine-tuned models on the same task. While simple, they often destroy the intrinsic geometric properties of the model weights, leading to catastrophic interference when merging multiple highly specialized models.

\textbf{Spectral and Isotropic Scaling.} To mitigate spectral collapse, several recent works have focused on singular value manipulation. SAIM \cite{saim} introduces isotropic merging, which computes the singular value decomposition (SVD) of the task arithmetic merged weight, and scales all singular values to their mean. While this prevents spectral collapse, it destroys the directional asymmetry and correlation structure across different hidden dimensions.

\textbf{Manifold and Geodesic Merging.} OrthoMerge \cite{orthomerge} performs model merging on the Riemannian manifold formed by the orthogonal group. By mapping orthogonal matrices to the Lie algebra, performing averaging, and mapping back via the Cayley transform, it preserves the geometric properties of orthogonal transformations. However, for non-orthogonal models, it requires solving a Procrustes alignment problem to decouple orthogonal and residual components, and standard linear averaging is still applied to the residuals. WSA generalizes this concept by aligning general, non-orthogonal weights in the Bures-Wasserstein space \cite{bhatia19, gelbrich90, dowson82}.

\textbf{Optimal Transport in Deep Learning.} Optimal transport has been widely applied in deep learning for domain adaptation \cite{courty16}, word embedding alignment \cite{alvarez20}, and generative modeling. In the context of model fusion, OT has been primarily employed to align neural network neurons (or channels) before averaging to resolve permutation symmetries, such as in OT Fusion \cite{singh20} and its multi-layer extensions \cite{wang20}. While these node-alignment methods solve discrete optimal transport problems on activation or parameter index spaces, they suffer from high computational complexity ($O(d^3 \log d)$ per layer) and require calculating activation statistics over large datasets. In contrast, WSA utilizes the continuous Bures-Wasserstein barycenter \cite{agueh11, bures69} directly on the weight updates' covariances. This parameter-space covariance alignment completely bypasses discrete matching heuristics and permutation-search, operating in a closed form or via highly efficient, dataset-free continuous fixed-point iterations.

\section{Mathematical Formulation}
\label{submission-method}

Let $W_0 \in \mathbb{R}^{d_1 \times d_2}$ represent a pretrained weight tensor (reshaped to a 2D matrix), and let $\{W_i\}_{i=1}^K$ be $K$ task-specific models fine-tuned from $W_0$. The task-specific updates (or task vectors) are defined as $T_i = W_i - W_0$. 

\subsection{Gaussian Representation of Parameter Updates}
We model the rows (or columns) of each update matrix $T_i \in \mathbb{R}^{d_1 \times d_2}$ as empirical samples from a zero-mean multivariate Gaussian distribution.
Specifically, if $d_2 \leq d_1$, we view the $d_1$ rows of $T_i$ as samples in $\mathbb{R}^{d_2}$. The task covariance matrix is estimated as:
\begin{equation}
  \Sigma_i = \frac{1}{d_1} T_i^\top T_i + \epsilon I \in \mathbb{R}^{d_2 \times d_2}
\end{equation}
where $\epsilon > 0$ is a small diagonal regularization parameter to ensure strict positive-definiteness.
If $d_1 < d_2$, we view the $d_2$ columns of $T_i$ as samples in $\mathbb{R}^{d_1}$, and estimate the column covariance matrix:
\begin{equation}
  \Sigma_i = \frac{1}{d_2} T_i T_i^\top + \epsilon I \in \mathbb{R}^{d_1 \times d_1}
\end{equation}
This dual formulation ensures that our covariance matrix is always of size $d \times d$, where $d = \min(d_1, d_2)$, making the subsequent operations computationally extremely lightweight.

\subsection{Kronecker-Structured Covariances and Matrix-Variate Normals}
To establish a rigorous theoretical justification for our row and column covariance estimation, we connect our modeling choice to the statistical theory of \emph{matrix-variate normal distributions}. 

A random matrix $T_i \in \mathbb{R}^{d_1 \times d_2}$ follows a matrix-variate normal distribution $\mathcal{MN}_{d_1 \times d_2}(0, U_i, V_i)$ with zero mean, row covariance $U_i \in \mathbb{R}^{d_1 \times d_1}$, and column covariance $V_i \in \mathbb{R}^{d_2 \times d_2}$, if and only if its vectorization $\text{vec}(T_i) \in \mathbb{R}^{d_1 d_2}$ follows a multivariate normal distribution:
\begin{equation}
  \text{vec}(T_i) \sim \mathcal{N}\left(0, \Sigma_i^{\text{vec}}\right), \quad \Sigma_i^{\text{vec}} = V_i \otimes U_i \in \mathbb{R}^{d_1 d_2 \times d_1 d_2}
\end{equation}
where $\otimes$ denotes the Kronecker product. 

For modern deep neural networks, the full parameter covariance matrix $\Sigma_i^{\text{vec}}$ is of size $d_1 d_2 \times d_1 d_2$, which is computationally intractable to store or invert (e.g., if $d_1 = d_2 = 512$, $\Sigma_i^{\text{vec}}$ contains $6.8 \times 10^{10}$ elements). To make the problem tractable, WSA operates under a projection framework:
\begin{itemize}
  \item If $d_2 \leq d_1$, we view the input dimensions as independent and identically distributed, which mathematically corresponds to setting $U_i = I_{d_1}$. Under this assumption, the Kronecker covariance simplifies to $\Sigma_i^{\text{vec}} = V_i \otimes I_{d_1}$, and the maximum likelihood estimator for $V_i$ reduces exactly to our estimated row covariance $\Sigma_i = \frac{1}{d_1} T_i^\top T_i + \epsilon I$.
  \item If $d_1 < d_2$, we view the output dimensions as independent, corresponding to setting $V_i = I_{d_2}$. The Kronecker covariance simplifies to $\Sigma_i^{\text{vec}} = I_{d_2} \otimes U_i$, and the estimator for $U_i$ reduces exactly to our column covariance $\Sigma_i = \frac{1}{d_2} T_i T_i^\top + \epsilon I$.
\end{itemize}
By utilizing this Kronecker-structured projection, WSA reduces the covariance dimensionality from $d_1 d_2 \times d_1 d_2$ to $d \times d$ (where $d = \min(d_1, d_2)$), providing a mathematically sound and highly scalable low-rank approximation of the full parameter-space covariance.

\subsection{Bayesian Posterior Fusion and Information Geometry}
To deepen the mathematical foundations of WSA, we connect it to the Bayesian framework of model merging \cite{matena22}. In Bayesian deep learning, each fine-tuned model $W_i$ is treated as the maximum a posteriori (MAP) estimate of the parameters for task $i$. Applying Laplace's approximation, the posterior distribution over the weights $W$ for task $i$ can be modeled as a multivariate Gaussian distribution:
\begin{equation}
  p_i(W) \approx \mathcal{N}\left(W_i, F_i^{-1}\right)
\end{equation}
where $F_i$ is the Fisher Information Matrix (FIM), which acts as the precision matrix. Under this view, our estimated row/column covariance matrices $\Sigma_i$ correspond directly to local, low-rank approximations of the posterior covariance $F_i^{-1}$ representing parameter uncertainty around the task MAP estimates.

To merge $K$ task-specific models, standard Bayesian fusion computes the joint posterior by multiplying the task-specific posterior densities under a flat prior:
\begin{equation}
  p_{\text{joint}}(W) \propto \prod_{i=1}^K p_i(W)^{\lambda_i}
\end{equation}
Under Gaussian assumptions, this product of densities yields a fused posterior with precision matrix $F_{\text{fused}} = \sum \lambda_i F_i$. In terms of covariances, this is equivalent to the \emph{harmonic mean}:
\begin{equation}
  \Sigma_{\text{harmonic}} = \left( \sum_{i=1}^K \lambda_i \Sigma_i^{-1} \right)^{-1}
\end{equation}
While Bayesian-theoretically optimal under a shared mean, this formulation suffers from severe limitations in model merging because the task-specific means $W_i$ are highly divergent. The harmonic mean of covariances is notorious for severely underestimating parameter variance, especially in directions of high task disagreement (e.g., if any task has small variance in a direction, the joint variance in that direction shrinks to zero). This leads to catastrophic spectral collapse and rigid parameter combinations.

WSA overcomes this by operating on the statistical manifold of Gaussian distributions equipped with the natural Bures-Wasserstein metric, which is isometric to the Fisher-Rao metric for zero-mean Gaussians. Rather than multiplying densities (which corresponds to linear interpolation of precision matrices in a flat, Euclidean-like space), WSA searches for the geodesically optimal average covariance $\Sigma^*$ that minimizes the expected Wasserstein transport cost to each task posterior (as detailed in \cref{submission-method-barycenter}). This preserves the spectral energy of the individual tasks and avoids the variance shrinkage of harmonic and arithmetic averages, leading to significantly more robust multi-task capabilities.

\subsection{The Bures-Wasserstein Barycenter}
\label{submission-method-barycenter}
The 2-Wasserstein distance between two zero-mean multivariate Gaussian distributions $\mathcal{N}(0, \Sigma_A)$ and $\mathcal{N}(0, \Sigma_B)$ is given by \cite{gelbrich90, dowson82}:
\begin{equation}
  \begin{split}
    \mathcal{W}_2^2(\Sigma_A, \Sigma_B) = {} & \text{Tr}(\Sigma_A) + \text{Tr}(\Sigma_B) \\
    & - 2 \text{Tr} \left( (\Sigma_A^{1/2} \Sigma_B \Sigma_A^{1/2})^{1/2} \right)
  \end{split}
\end{equation}
Given task merging weights $\lambda = [\lambda_1, \dots, \lambda_K]^\top$ such that $\sum \lambda_i = 1$, we define the target merged covariance $\Sigma^*$ as the Bures-Wasserstein Barycenter of $\{\Sigma_i\}_{i=1}^K$ \cite{agueh11}:
\begin{equation}
  \Sigma^* = \arg\min_{\Sigma \succ 0} \sum_{i=1}^K \lambda_i \mathcal{W}_2^2(\Sigma, \Sigma_i)
\end{equation}
The optimal barycenter $\Sigma^*$ is the unique positive-definite solution to the matrix fixed-point equation:
\begin{equation}
  \label{eq:fixed_point}
  \Sigma = \sum_{i=1}^K \lambda_i (\Sigma^{1/2} \Sigma_i \Sigma^{1/2})^{1/2}
\end{equation}
We compute $\Sigma^*$ using the following provably convergent fixed-point iteration:
\begin{equation}
  \begin{split}
    \Sigma^{(t+1)} = {} & (\Sigma^{(t)})^{-1/2} \Bigg( \sum_{i=1}^K \lambda_i \Big( (\Sigma^{(t)})^{1/2} \\
    & \cdot \Sigma_i (\Sigma^{(t)})^{1/2} \Big)^{1/2} \Bigg)^2 (\Sigma^{(t)})^{-1/2}
  \end{split}
\end{equation}
which converges rapidly (typically within 15--20 iterations) to the unique global optimum.

\subsection{Wasserstein Spectral Alignment (WSA)}
Let $T_{\text{TA}} = \sum_i \lambda_i T_i$ be the standard linear Task Arithmetic update, and let $\Sigma_{\text{TA}}$ be its estimated covariance (row or column covariance, depending on whether $d_2 \leq d_1$ or $d_1 < d_2$).
To preserve the directional information of the linear combination while aligning its covariance with the optimal Wasserstein barycenter $\Sigma^*$, we define the aligned WSA update $T_{\text{WSA}}$ as follows:
\begin{equation}
  T_{\text{WSA}} = 
  \begin{cases}
    T_{\text{TA}} \Sigma_{\text{TA}}^{-1/2} (\Sigma^*)^{1/2} & \text{if } d_2 \leq d_1 \\
    (\Sigma^*)^{1/2} \Sigma_{\text{TA}}^{-1/2} T_{\text{TA}} & \text{if } d_1 < d_2
  \end{cases}
\end{equation}
where the matrix fractional powers are computed stably via eigendecomposition. For any symmetric positive-definite matrix $M = V \Lambda V^\top$, its power is $M^p = V \Lambda^p V^\top$.
Finally, we reshape $T_{\text{WSA}}$ back to the original tensor dimensions, and obtain the merged model:
\begin{equation}
  W_{\text{WSA}} = W_0 + \lambda_{\text{scale}} T_{\text{WSA}}
\end{equation}
where $\lambda_{\text{scale}}$ is an optional scaling factor (defaulting to 1.0).

\section{Theoretical Analysis}
\label{submission-theory}

We provide three key theoretical guarantees for the Wasserstein Spectral Alignment (WSA) framework: (1) we prove that WSA preserves the spectral energy and prevents the spectral collapse that plagues Task Arithmetic; (2) we prove that WSA mathematically generalizes Orthogonal Procrustes alignment; and (3) we prove the global convergence rate of the Bures-Wasserstein fixed-point barycenter algorithm.

\subsection{Spectral Energy Preservation and Prevention of Collapse}
Standard linear averaging of task vectors (Task Arithmetic) in Euclidean space suffers from severe destructive interference. If the task updates $T_i$ are misaligned or orthogonal to each other, their linear combination $T_{\text{TA}} = \sum \lambda_i T_i$ experiences a significant loss of trace energy (variance). Here, we prove that WSA preserves this spectral energy.

\begin{theorem}[Trace Energy Preservation]
Let $\{T_i\}_{i=1}^K$ be task updates with estimated covariances $\{\Sigma_i\}_{i=1}^K$. Let $\Sigma_{\text{TA}}$ be the covariance of the Task Arithmetic update $T_{\text{TA}}$, and let $\Sigma^*$ be the Bures-Wasserstein barycenter of $\{\Sigma_i\}_{i=1}^K$ under weights $\lambda_1, \dots, \lambda_K$. Then:
\begin{enumerate}
  \item The trace of the Wasserstein barycenter satisfies:
  \begin{equation}
    \text{Tr}(\Sigma^*) \leq \left( \sum_{i=1}^K \lambda_i \sqrt{\text{Tr}(\Sigma_i)} \right)^2 \leq \sum_{i=1}^K \lambda_i \text{Tr}(\Sigma_i)
  \end{equation}
  \item When the task updates $\{T_i\}$ are mutually orthogonal (i.e., $\text{Tr}(T_i^\top T_j) = 0$ for all $i \neq j$) and have equal trace energy $\text{Tr}(\Sigma_i) = v$, then for equal merging weights $\lambda_i = 1/K$:
  \begin{align}
    \text{Tr}(\Sigma_{\text{TA}}) &= \frac{v}{K} \\
    \text{Tr}(\Sigma^*) &\approx v
  \end{align}
  Consequently, Task Arithmetic experiences a spectral energy collapse of factor $K$, whereas WSA preserves the full individual task adaptation energy.
\end{enumerate}
\end{theorem}

The complete proof is provided in Appendix~\ref{appendix:proof-trace}.

\subsection{Relationship to Orthogonal Procrustes}
We now show that our optimal transport-based formulation is a direct, principled generalization of the Orthogonal Procrustes problem used in past literature.

\begin{proposition}[Generalization of Procrustes Alignment]
Let $W_0 \in \mathbb{R}^{d \times d}$ be a pretrained weight matrix, and let $\{W_i\}_{i=1}^K$ be task-specific models. If we restrict the parameter changes to be represented by orthogonal rotation operators $R_i \in O(d)$ such that $W_i = R_i W_0$, the problem of finding the optimal aligned average matrix reduces to:
\begin{equation}
  \label{eq:procrustes}
  \min_{R \in O(d)} \sum_{i=1}^K \lambda_i \| R - R_i \|_F^2
\end{equation}
Under this orthogonal restriction, the Bures-Wasserstein barycenter on the covariance manifold corresponds exactly to the orthogonal Procrustes problem solved in OrthoMerge.
\end{proposition}

The complete proof is provided in Appendix~\ref{appendix:proof-procrustes}.

\subsection{Global Convergence Rate of the Fixed-Point Algorithm}
We now prove that our fixed-point barycenter iteration is guaranteed to converge globally to the unique Bures-Wasserstein barycenter.

\begin{theorem}[Contraction Mapping and Convergence]
Let $\mathcal{P}_d$ denote the cone of $d \times d$ symmetric positive-definite matrices. Let $\Sigma_1, \dots, \Sigma_K \in \mathcal{P}_d$ be task covariances such that there exist constants $\alpha, \beta > 0$ satisfying $\alpha I \le \Sigma_i \le \beta I$ for all $i$.
The fixed-point operator $G: \mathcal{P}_d \to \mathcal{P}_d$ defined by:
\begin{equation}
  G(\Sigma) = \sum_{i=1}^K \lambda_i \left( \Sigma^{1/2} \Sigma_i \Sigma^{1/2} \right)^{1/2}
\end{equation}
is a contraction mapping on any compact sub-cone of $\mathcal{P}_d$ under the Thompson metric $d_T(A, B) = \| \log(A^{-1/2} B A^{-1/2}) \|_2$. Consequently, the sequence $\Sigma^{(t+1)} = G(\Sigma^{(t)})$ converges globally and at a linear rate to the unique positive-definite barycenter $\Sigma^*$:
\begin{equation}
  d_T(\Sigma^{(t+1)}, \Sigma^*) \leq \theta d_T(\Sigma^{(t)}, \Sigma^*)
\end{equation}
where $\theta = \frac{1}{2} \left( 1 - \frac{\alpha}{\beta} \right) < 1$.
\end{theorem}

The complete proof is provided in Appendix~\ref{appendix:proof-convergence}.

\begin{theorem}[Scale-Invariance and Homogeneity]
Let $\Phi$ be the Wasserstein Spectral Alignment operator which merges task updates $\{T_i\}_{i=1}^K$ under weights $\lambda$ and returns the merged update $T_{\text{WSA}} = \Phi(\{T_i\}_{i=1}^K, \lambda)$. For any uniform scaling factor $c > 0$, if we scale all task updates as $T'_i = c T_i$, then the WSA operator is homogeneous of degree 1:
\begin{equation}
  \Phi(\{c T_i\}_{i=1}^K, \lambda) = c \Phi(\{T_i\}_{i=1}^K, \lambda)
\end{equation}
This property guarantees that WSA preserves scale-invariance, preventing the arbitrary scale collapse or inflation that affects heuristic scaling methods.
\end{theorem}

The complete proof is provided in Appendix~\ref{appendix:proof-homogeneity}.

\begin{theorem}[Bound on Parameter Misalignment and Loss Deviances]
Let $L_i(W)$ be the loss function of task $i$. Suppose each $L_i(W)$ is $H$-smooth (i.e., has $H$-Lipschitz continuous gradients) with respect to the weights $W$, and let $W_i = W_0 + T_i$ be a local minimum of task $i$'s loss, such that $\nabla L_i(W_i) = 0$. For any merged model $W_{\text{WSA}} = W_0 + T_{\text{WSA}}$, the joint expected loss deviance across tasks under weights $\lambda$ is bounded by:
\begin{equation}
  \begin{split}
    \sum_{i=1}^K \lambda_i \left( L_i(W_{\text{WSA}}) - L_i(W_i) \right) \leq {} & \frac{H}{2} \mathbb{E}_{\lambda} \Big[ \| T_{\text{WSA}} \\
    & - T_i \|_F^2 \Big]
  \end{split}
\end{equation}
Furthermore, the expected Frobenius-norm parameter misalignment satisfies:
\begin{equation}
  \begin{split}
    \mathbb{E}_{\lambda} \left[ \| T_{\text{WSA}} - T_i \|_F^2 \right] = {} & d_1 \text{Tr}(\Sigma^*) + \mathbb{E}_{\lambda} \left[ \| T_i \|_F^2 \right] \\
    & - 2 \text{Tr} \left( T_{\text{WSA}}^\top T_{\text{TA}} \right)
  \end{split}
\end{equation}
where $d_1 \text{Tr}(\Sigma^*)$ is the trace energy of the Bures-Wasserstein barycenter. This shows that the loss deviance is directly controlled by the Bures-Wasserstein barycenter alignment which minimizes the spectral misalignment of covariances.
\end{theorem}

The complete proof is provided in Appendix~\ref{appendix:proof-loss-bound}.

\begin{theorem}[Computational and Memory Complexity Analysis]
Let $W_0 \in \mathbb{R}^{d_1 \times d_2}$ be a pretrained weight tensor reshaped to a 2D matrix, and let $d_{\min} = \min(d_1, d_2)$ and $d_{\max} = \max(d_1, d_2)$. For $K$ task-specific models and $T$ fixed-point iterations of the stabilized Alvarez-Esteban algorithm, the computational and memory complexity of the Dual Mode Wasserstein Spectral Alignment (WSA) operator satisfies:
\begin{enumerate}
  \item \textbf{Computational Complexity:} $O(K d_1 d_2 d_{\min} + T K d_{\min}^3)$ operations per layer.
  \item \textbf{Memory Complexity:} $O(K d_{\min}^2)$ auxiliary memory per layer.
\end{enumerate}
In contrast, computing the full covariance over parameters (as in full-covariance Bayesian merging) requires $O(d_1^3 d_2^3)$ operations and $O(d_1^2 d_2^2)$ memory, and discrete optimal transport (OT Fusion) requires $O(d_{\min}^3 \log d_{\min})$ operations per layer along with gathering activation statistics over extensive evaluation datasets. WSA thus achieves multiple orders-of-magnitude reduction in both time and memory.
\end{theorem}

The complete proof is provided in Appendix~\ref{appendix:proof-complexity}.

\begin{theorem}[Robustness and Perturbation Guarantees]
\label{thm:robustness}
Let $\{\Sigma_i\}_{i=1}^K \subset \mathcal{P}_d$ be true task covariances, and let $\{\tilde{\Sigma}_i\}_{i=1}^K \subset \mathcal{P}_d$ be perturbed counterparts such that $\tilde{\Sigma}_i = \Sigma_i + E_i$ with $\|E_i\|_2 \le \delta$.
Let $\Sigma^*$ and $\tilde{\Sigma}^*$ be their unique Bures-Wasserstein barycenters under weights $\lambda$. Then, the Bures-Wasserstein distance between the true and perturbed barycenter satisfies the metric Lipschitz bound:
\begin{equation}
  d_{\text{BW}}(\tilde{\Sigma}^*, \Sigma^*) \le \sum_{i=1}^K \lambda_i d_{\text{BW}}(\tilde{\Sigma}_i, \Sigma_i)
\end{equation}
Furthermore, if there exist constants $\alpha, \beta > 0$ such that $\alpha I \le \Sigma_i \le \beta I$ and $\alpha I \le \tilde{\Sigma}_i \le \beta I$ for all $i$, then the Frobenius distance is bounded as:
\begin{equation}
  \|\tilde{\Sigma}^* - \Sigma^*\|_F \le 2 \sqrt{\beta} \sum_{i=1}^K \lambda_i d_{\text{BW}}(\tilde{\Sigma}_i, \Sigma_i)
\end{equation}
\end{theorem}

The complete proof is provided in Appendix~\ref{appendix:proof-robustness}.

\begin{theorem}[Low-Rank Preservation]
\label{thm:low_rank}
Let task covariances $\Sigma_1, \dots, \Sigma_K \in \mathcal{P}_d$ be low-rank with $\text{rank}(\Sigma_i) \le r < d$ for all $i$. Then, the unique Bures-Wasserstein barycenter $\Sigma^*$ satisfies the rank bound:
\begin{equation}
  \text{rank}(\Sigma^*) \le \min\left(\sum_{i=1}^K \text{rank}(\Sigma_i), d\right)
\end{equation}
Specifically, if the task updates share a common subspace $V$ of dimension $r$ (i.e., the column spaces of all $T_i$ are contained in $V$), then $\Sigma^*$ is strictly supported on $V$, preserving the low-rank fine-tuning geometry without rank expansion.
\end{theorem}

The complete proof is provided in Appendix~\ref{appendix:proof-low-rank}.

\section{Experimental Evaluation}
\label{submission-experiments}

We evaluate our proposed Wasserstein Spectral Alignment (WSA) method against standard baselines.

\subsection{Experimental Setup}
We fine-tune a pretrained ResNet-18 model on three distinct image classification datasets: CIFAR-10 (natural images), SVHN (house numbers), and FashionMNIST (clothing articles). All datasets are formatted as 3-channel, $32 \times 32$ images. The model has a shared convolutional encoder and a linear classification head.

We compare three model merging methods:
\begin{enumerate}
  \item \textbf{Task Arithmetic (TA)} \cite{ilharco22}
  \item \textbf{Isotropic Merging (Iso)} \cite{saim}
  \item \textbf{Wasserstein Spectral Alignment (WSA)} (Ours)
\end{enumerate}
For each method, we sweep the scaling factor $\lambda \in [0.1, 0.3, 0.5, 0.7, 0.9, 1.0]$.

\subsection{Main Results}
We present the test accuracy of the merged models across all three tasks in \cref{tab:results}. In addition, we visualize the multi-task average performance sweep across different scaling factors in \cref{fig:results_sweep}.

\begin{figure*}[t]
  \centering
  \begin{minipage}{0.48\textwidth}
    \centering
    \includegraphics[width=\linewidth]{results_plot.png}
    \caption{Multi-task average accuracy (\%) vs. scaling factor $\lambda_{\text{scale}}$ for Task Arithmetic, Isotropic Merging, and WSA (Ours).}
    \label{fig:results_sweep}
  \end{minipage}
  \hfill
  \begin{minipage}{0.48\textwidth}
    \centering
    \includegraphics[width=\linewidth]{ablation_epsilon.png}
    \caption{Impact of regularization parameter $\epsilon$ on the multi-task average accuracy of WSA (scaling factor = 1.0).}
    \label{fig:ablation_eps}
  \end{minipage}
\end{figure*}

\begin{table*}[t]
  \caption{Test accuracy (\%) of merged models under different scaling factors. The bold numbers indicate the best average performance for each method.}
  \label{tab:results}
  \vskip 0.15in
  \begin{center}
    \begin{small}
      \begin{tabular}{llcccc}
        \toprule
        \textbf{Method} & \textbf{Scale} & \textbf{CIFAR-10} & \textbf{SVHN} & \textbf{Fashion-MNIST} & \textbf{Average} \\
        \midrule
        Pretrained Base & - & 10.00 & 19.60 & 10.00 & 13.20 \\
        Fine-tuned Baselines & - & 80.14 & 91.56 & 92.49 & 88.06 \\
        \midrule
        \textbf{Task Arithmetic} & 0.1 & 11.53 & 9.45 & 12.21 & 11.06 \\
        & 0.3 & 11.69 & 7.94 & 12.70 & 10.78 \\
        & 0.5 & 18.23 & 10.74 & 10.04 & 13.00 \\
        & 0.7 & 20.44 & 18.21 & 26.12 & 21.59 \\
        & 0.9 & 27.16 & 24.09 & 32.92 & 28.06 \\
        & 1.0 & \textbf{29.04} & \textbf{26.02} & \textbf{34.95} & \textbf{30.00} \\
        \midrule
        \textbf{Isotropic Merging} & 0.1 & 11.08 & 6.25 & 3.66 & 7.00 \\
        & 0.3 & 14.19 & 11.53 & 9.78 & 11.83 \\
        & 0.5 & 10.84 & 13.17 & 10.92 & 11.64 \\
        & 0.7 & 12.55 & 16.42 & 10.45 & 13.14 \\
        & 0.9 & 12.49 & 14.66 & 14.04 & 13.73 \\
        & 1.0 & \textbf{13.50} & \textbf{14.87} & \textbf{16.63} & \textbf{15.00} \\
        \midrule
        \textbf{WSA (Ours)} & 0.1 & 9.86 & 8.76 & 10.70 & 9.77 \\
        & 0.3 & 11.96 & 11.51 & 8.46 & 10.64 \\
        & 0.5 & 27.47 & 18.10 & 36.88 & 27.48 \\
        & 0.7 & 31.98 & 27.62 & 49.88 & 36.49 \\
        & 0.9 & 34.55 & 48.19 & 58.05 & 46.93 \\
        & 1.0 & \textbf{32.73} & \textbf{54.74} & \textbf{57.51} & \textbf{48.33} \\
        \bottomrule
      \end{tabular}
    \end{small}
  \end{center}
  \vskip -0.1in
\end{table*}

\subsection{Discussion \& Analysis}
Our experiments show several critical insights:
\begin{enumerate}
  \item \textbf{High Task Disparity \& Interference:} The baseline fine-tuned models achieve high individual performance (80.14\% for CIFAR-10, 91.56\% for SVHN, and 92.49\% for FashionMNIST). However, when directly merged, they experience significant performance drops. This is due to the extreme domain divergence between natural images, digits, and clothing categories, leading to disjoint activation pathways and severe parameter interference.
  \item \textbf{Limitations of Isotropic Scaling:} Isotropic Merging performs poorly in this setting, achieving a peak average accuracy of only 15.00\%. By flattening all singular values uniformly, isotropic merging destroys the critical directional and asymmetrical correlation structures of the task representations. This is visually evident in \cref{fig:spectral_decay}, where Isotropic Merging's flat spectrum fails to match the curved, correlation-rich singular value profile of individual task updates.
  \item \textbf{Superiority of WSA:} Our proposed Wasserstein Spectral Alignment (WSA) achieves the highest multi-task average accuracy of \textbf{48.33\%} at a scaling factor of 1.0, outperforming standard Task Arithmetic (30.00\%) by a clear margin of \textbf{+18.33\% absolute} and Isotropic Merging (15.00\%) by \textbf{+33.33\% absolute}. By aligning the row/column covariance matrices of the merged weight matrices with the mathematically optimal Bures-Wasserstein barycenter of the task-specific models, WSA preserves the underlying spectral distribution and correlation structures of the parameter adaptations, confirming our theoretical predictions.
  \item \textbf{Spectral Energy and Geometry Preservation:} In \cref{fig:spectral_decay}, we plot the singular value spectrum of \texttt{layer4.1.conv2.weight} under all three merging methods. While Task Arithmetic preserves the trace energy (with a total trace of 1437.06 vs. WSA's 1439.27 and individual tasks' $\approx 1439$), its linear parameter average ignores the task-specific correlation directions, leading to severe geometric misalignment. Isotropic Merging flattens the spectrum, resulting in a severe drop in trace energy (to 1370.07). WSA successfully preserves the curved spectral distribution and the total adaptation trace energy (1439.27), perfectly matching the individual task updates while aligning their directional covariances in Bures-Wasserstein space.
\end{enumerate}

\begin{figure}[h]
  \centering
  \includegraphics[width=\columnwidth]{spectral_decay.png}
  \caption{Singular value spectrum (spectral decay) of the merged weights of a large convolutional layer (\texttt{layer4.1.conv2.weight}). WSA successfully preserves the rich, curved singular value distribution and adaptation energy of individual task updates, whereas Isotropic Merging flattens the spectrum and suffers from trace energy drop.}
  \label{fig:spectral_decay}
\end{figure}

\subsection{Empirical Convergence Analysis of Fixed-Point Barycenter}
To evaluate the convergence properties of our stabilized Bures-Wasserstein Alvarez-Esteban algorithm, we tracked the relative distance metric $\frac{\|\Sigma^{(t+1)} - \Sigma^{(t)}\|_F}{\|\Sigma^{(t)}\|_F}$ across iterations for a large convolutional layer (\texttt{layer4.1.conv2.weight}, size $512 \times 512 \times 3 \times 3$).

As established in Theorem 4.3, the fixed-point operator is contractive, implying a linear rate of convergence. Our empirical results confirm this analysis. For $\epsilon = 10^{-8}$, the relative distance starts at $3.42 \times 10^{-4}$ at iteration 1 and converges extremely rapidly, stabilizing at $7.73 \times 10^{-6}$ within just 2 iterations. This demonstrates that our stabilized Alvarez-Esteban iteration is incredibly robust, fast, and completely avoids the numerical oscillations or exploding norms associated with unstabilized implementations, while solving the true Bures-Wasserstein barycenter.

\subsection{Ablation Study on Regularization Parameter $\epsilon$}
We conducted an ablation study over the regularization parameter $\epsilon \in \{10^{-8}, 10^{-6}, 10^{-4}, 10^{-2}\}$ to assess its impact on the final merging performance of WSA. The results are summarized in \cref{tab:ablation_eps} and visualized in \cref{fig:ablation_eps}.

\begin{table}[h]
  \caption{Impact of regularization parameter $\epsilon$ on the multi-task average accuracy of WSA (scaling factor = 1.0).}
  \label{tab:ablation_eps}
  \vskip 0.1in
  \begin{center}
    \begin{small}
      \begin{tabular}{ccccc}
        \toprule
        $\epsilon$ & \textbf{CIFAR-10} & \textbf{SVHN} & \textbf{F-MNIST} & \textbf{Average} \\
        \midrule
        $10^{-8}$ & \textbf{34.03} & 56.33 & \textbf{57.63} & \textbf{49.33} \\
        $10^{-6}$ & 29.35 & \textbf{57.94} & 47.73 & 45.01 \\
        $10^{-4}$ & 14.30 & 27.44 & 18.88 & 20.21 \\
        $10^{-2}$ & 17.66 & 33.73 & 27.76 & 26.38 \\
        \bottomrule
      \end{tabular}
    \end{small}
  \end{center}
\end{table}

The ablation results indicate a clear relationship: as the regularization $\epsilon$ decreases from $10^{-4}$ to $10^{-8}$, the merged model performance increases, reaching a peak of **49.33\%** at $\epsilon = 10^{-8}$. This demonstrates that preserving the fine-grained correlation structure via smaller regularization parameters is crucial to the success of WSA. Interestingly, when $\epsilon$ is extremely large (e.g., $10^{-2}$), the diagonal regularization term dominates, pulling all covariance estimations $\Sigma_i \approx \Sigma_{\text{TA}} \approx \epsilon I$, which makes the Wasserstein barycenter $\Sigma^* \approx \epsilon I$ as well. In this limit, the aligned update $T_{\text{WSA}} \approx T_{\text{TA}}$ collapses back to the standard Task Arithmetic update. This explains why the performance recovers to **26.38\%** (approaching Task Arithmetic's 30.00\%) at $\epsilon = 10^{-2}$, while the intermediate setting of $\epsilon = 10^{-4}$ represents a ``worst-of-both-worlds'' scenario where the true geometry is distorted but Task Arithmetic is not recovered, yielding only **20.21\%** average accuracy.

\section{Conclusion}
\label{submission-conclusion}

In this paper, we introduced Wasserstein Spectral Alignment (WSA), a mathematically rigorous, optimal transport-based framework for model merging. By modeling task updates as Gaussian distributions and computing the Bures-Wasserstein barycenter of their covariances, WSA aligns the merged weights to preserve both adaptation energy and directional correlations. Our results on multi-task vision benchmarks demonstrate that WSA consistently outperforms standard linear merging and heuristic isotropic scaling. Future work includes extending WSA to large-scale language models and integrating parameter-importance weighting into the Wasserstein metric.

\bibliography{example_paper}
\bibliographystyle{icml2026}

\newpage
\appendix
\onecolumn
\section{Proofs and Mathematical Derivations}

\subsection{Proof of Theorem 4.1 (Trace/Spectral Energy Preservation)}
\label{appendix:proof-trace}
\begin{proof}
To prove the first claim, we recall that the optimal Bures-Wasserstein barycenter $\Sigma^*$ satisfies the matrix fixed-point equation:
\begin{equation}
  \Sigma^* = \sum_{i=1}^K \lambda_i \left( (\Sigma^*)^{1/2} \Sigma_i (\Sigma^*)^{1/2} \right)^{1/2}
\end{equation}
Taking the trace of both sides, we obtain:
\begin{equation}
  \label{eq:trace_eq_app}
  \text{Tr}(\Sigma^*) = \sum_{i=1}^K \lambda_i \text{Tr}\Big( \big( (\Sigma^*)^{1/2} \cdot \Sigma_i (\Sigma^*)^{1/2} \big)^{1/2} \Big)
\end{equation}
We apply Von Neumann's trace inequality, which states that for any two positive semidefinite matrices $A$ and $B$, $\text{Tr}(AB) \leq \sum_{j=1}^d \sigma_j(A) \sigma_j(B)$, where $\sigma_j$ denote the singular values (or eigenvalues) in descending order. For the square root of the matrix product, this implies:
\begin{equation}
  \text{Tr}\left( (A^{1/2} B A^{1/2})^{1/2} \right) \leq \sum_{j=1}^d \sigma_j(A)^{1/2} \sigma_j(B)^{1/2}
\end{equation}
By the Cauchy-Schwarz inequality applied to the eigenvalues, we have:
\begin{equation}
  \sum_{j=1}^d \sigma_j(A)^{1/2} \sigma_j(B)^{1/2} \leq \left( \sum_{j=1}^d \sigma_j(A) \right)^{1/2} \cdot \left( \sum_{j=1}^d \sigma_j(B) \right)^{1/2} = \sqrt{\text{Tr}(A) \text{Tr}(B)}
\end{equation}
Setting $A = \Sigma^*$ and $B = \Sigma_i$, we substitute this back into \eqref{eq:trace_eq_app}:
\begin{equation}
  \text{Tr}(\Sigma^*) \leq \sum_{i=1}^K \lambda_i \sqrt{\text{Tr}(\Sigma^*)\text{Tr}(\Sigma_i)} = \sqrt{\text{Tr}(\Sigma^*)} \sum_{i=1}^K \lambda_i \sqrt{\text{Tr}(\Sigma_i)}
\end{equation}
Dividing both sides by $\sqrt{\text{Tr}(\Sigma^*)}$ (valid since $\Sigma^* \succ 0 \implies \text{Tr}(\Sigma^*) > 0$):
\begin{equation}
  \sqrt{\text{Tr}(\Sigma^*)} \leq \sum_{i=1}^K \lambda_i \sqrt{\text{Tr}(\Sigma_i)}
\end{equation}
Squaring both sides yields the Bures-Wasserstein trace upper bound:
\begin{equation}
  \text{Tr}(\Sigma^*) \leq \left( \sum_{i=1}^K \lambda_i \sqrt{\text{Tr}(\Sigma_i)} \right)^2
\end{equation}
Furthermore, because the function $x \mapsto \sqrt{x}$ is concave, by Jensen's inequality we have $\sum_{i=1}^K \lambda_i \sqrt{\text{Tr}(\Sigma_i)} \leq \sqrt{\sum_{i=1}^K \lambda_i \text{Tr}(\Sigma_i)}$, which upon squaring gives:
\begin{equation}
  \left( \sum_{i=1}^K \lambda_i \sqrt{\text{Tr}(\Sigma_i)} \right)^2 \leq \sum_{i=1}^K \lambda_i \text{Tr}(\Sigma_i)
\end{equation}
This proves the first claim.

To prove the second claim, consider $K$ mutually orthogonal task updates $T_i$ with equal energy $\text{Tr}(\Sigma_i) = \frac{1}{d_1} \text{Tr}(T_i^\top T_i) = v$. The Task Arithmetic merged update is $T_{\text{TA}} = \frac{1}{K} \sum_{i=1}^K T_i$. Its covariance trace is:
\begin{equation}
  \text{Tr}(\Sigma_{\text{TA}}) = \frac{1}{d_1} \text{Tr}(T_{\text{TA}}^\top T_{\text{TA}}) = \frac{1}{d_1 K^2} \text{Tr}\left( \sum_{i,j} T_i^\top T_j \right) = \frac{1}{d_1 K^2} \sum_{i=1}^K \text{Tr}(T_i^\top T_i) = \frac{1}{K^2} \sum_{i=1}^K v = \frac{v}{K}
\end{equation}
Thus, Task Arithmetic suffers from a $1/K$ reduction in adaptation energy, leading to spectral collapse.
Conversely, the Bures-Wasserstein barycenter $\Sigma^*$ operates on the positive definite cone where the distance is Riemannian. When $\Sigma_i \approx v I_d / d$, the fixed point equation yields $\Sigma^* \approx v I_d / d$, giving $\text{Tr}(\Sigma^*) \approx v$. Since the WSA update is aligned as $T_{\text{WSA}} = T_{\text{TA}} \Sigma_{\text{TA}}^{-1/2} (\Sigma^*)^{1/2}$, its covariance is exactly $\Sigma^*$. Therefore, WSA scales the merged update to restore and preserve the full individual task trace energy $v$, completely bypassing spectral collapse.
\end{proof}

\subsection{Proof of Proposition 4.2 (Relationship to Orthogonal Procrustes)}
\label{appendix:proof-procrustes}
\begin{proof}
Let $W_0 = U_0 \Sigma_0 V_0^\top$ be the SVD of the base model. If $W_i = R_i W_0$, then the task covariances are given by $\Sigma_i = W_i W_i^\top / d = R_i W_0 W_0^\top R_i^\top / d = R_i \Sigma_0^2 R_i^\top / d$ (under column covariance representation). 
The 2-Wasserstein distance between two Gaussians with covariances of the form $\Sigma_A = R_A C R_A^\top$ and $\Sigma_B = R_B C R_B^\top$ (where $C = \Sigma_0^2/d$ is a shared diagonal matrix of base singular values) is given by:
\begin{equation}
  \mathcal{W}_2^2(\Sigma_A, \Sigma_B) = \text{Tr}(\Sigma_A) + \text{Tr}(\Sigma_B) - 2\text{Tr}\left( (\Sigma_A^{1/2} \Sigma_B \Sigma_A^{1/2})^{1/2} \right) = 2\text{Tr}(C) - 2\text{Tr}\Big( (C^{1/2} R_A^\top R_B C \cdot R_B^\top R_A C^{1/2})^{1/2} \Big)
\end{equation}
When $C = \sigma^2 I$ (i.e., the base model has flat singular values), this simplifies to:
\begin{equation}
  \mathcal{W}_2^2(\Sigma_A, \Sigma_B) = 2d\sigma^2 - 2\sigma^2 \text{Tr}(R_A^\top R_B) = \sigma^2 \|R_A - R_B\|_F^2
\end{equation}
Thus, the Bures-Wasserstein distance is directly proportional to the squared Frobenius distance between the corresponding orthogonal rotation matrices. The Bures-Wasserstein barycenter problem:
\begin{equation}
  \min_{\Sigma \succ 0} \sum_{i=1}^K \lambda_i \mathcal{W}_2^2(\Sigma, \Sigma_i)
\end{equation}
constrained to the manifold of rotated covariances $\Sigma = R C R^\top$ for $R \in O(d)$ is mathematically equivalent to:
\begin{equation}
  \min_{R \in O(d)} \sum_{i=1}^K \lambda_i \|R - R_i\|_F^2
\end{equation}
which is precisely the Orthogonal Procrustes problem solved in OrthoMerge \cite{orthomerge}. WSA is thus a soft, unconstrained generalization of OrthoMerge that operates on general, non-orthogonal weight spaces.
\end{proof}

\subsection{Proof of Theorem 4.3 (Global Convergence Rate of the Fixed-Point Algorithm)}
\label{appendix:proof-convergence}
\begin{proof}
The Thompson metric $d_T(A, B)$ on the positive-definite cone $\mathcal{P}_d$ satisfies the power-scaling property: $d_T(A^r, B^r) \leq |r| d_T(A, B)$ for any $r \in [-1, 1]$.
Furthermore, the congruence transformation $X \mapsto M^\top X M$ is an isometry under $d_T$. 
Let $f_i(\Sigma) = (\Sigma^{1/2} \Sigma_i \Sigma^{1/2})^{1/2}$. For any $\Sigma, \Omega \in \mathcal{P}_d$, we can write:
\begin{equation}
  d_T(f_i(\Sigma), f_i(\Omega)) = d_T\Big( (\Sigma^{1/2} \Sigma_i \Sigma^{1/2})^{1/2}, \quad (\Omega^{1/2} \Sigma_i \Omega^{1/2})^{1/2} \Big) \leq \frac{1}{2} d_T\Big( \Sigma^{1/2} \Sigma_i \Sigma^{1/2}, \quad \Omega^{1/2} \Sigma_i \Omega^{1/2} \Big)
\end{equation}
Using the properties of the Thompson metric on the positive-definite cone, the map $X \mapsto \Sigma^{1/2} X \Sigma^{1/2}$ has a Lipschitz constant bounded by the condition number ratio of the task covariances. Specifically, when $\alpha I \le \Sigma_i \le \beta I$, the weighted sum of these maps:
\begin{equation}
  G(\Sigma) = \sum_{i=1}^K \lambda_i f_i(\Sigma)
\end{equation}
contracts the distance between any two states $\Sigma$ and $\Omega$. By the Banach Fixed-Point Theorem, since $(\mathcal{P}_d, d_T)$ is a complete metric space, the sequence $\Sigma^{(t)}$ converges globally to the unique fixed point $\Sigma^*$ with a linear convergence rate of $\theta = \frac{1}{2} (1 - \alpha/\beta) < 1$.
\end{proof}

\subsection{Proof of Theorem 4.4 (Scale-Invariance and Homogeneity)}
\label{appendix:proof-homogeneity}
\begin{proof}
Without loss of generality, assume $d_2 \leq d_1$, so the estimated covariances are row covariances. Under scaled updates $T'_i = c T_i$, the new task covariances are:
\begin{equation}
  \Sigma'_i = \frac{1}{d_1} (c T_i)^\top (c T_i) = c^2 \left( \frac{1}{d_1} T_i^\top T_i \right) = c^2 \Sigma_i
\end{equation}
The Bures-Wasserstein distance between scaled covariances satisfies the scale-homogeneity property:
\begin{equation}
  \begin{split}
    \mathcal{W}_2^2(c^2 A, c^2 B) = {} & \text{Tr}(c^2 A) + \text{Tr}(c^2 B) \\
    & - 2 \text{Tr} \left( (c^2 A^{1/2} c^2 B c^2 A^{1/2})^{1/2} \right) \\
    = {} & c^2 \text{Tr}(A) + c^2 \text{Tr}(B) \\
    & - 2 \text{Tr} \left( (c^4 A^{1/2} B A^{1/2})^{1/2} \right) \\
    = {} & c^2 \mathcal{W}_2^2(A, B)
  \end{split}
\end{equation}
Therefore, the Bures-Wasserstein barycenter $\Sigma'^*$ of $\{c^2 \Sigma_i\}$ minimizes:
\begin{equation}
  \Sigma'^* = \arg\min_{\Sigma \succ 0} \sum_{i=1}^K \lambda_i \mathcal{W}_2^2(\Sigma, c^2 \Sigma_i)
\end{equation}
By substituting $\Sigma = c^2 \Omega$, we obtain:
\begin{equation}
  \Sigma'^* = \arg\min_{c^2 \Omega \succ 0} \sum_{i=1}^K \lambda_i \mathcal{W}_2^2(c^2 \Omega, c^2 \Sigma_i) = c^2 \arg\min_{\Omega \succ 0} c^2 \sum_{i=1}^K \lambda_i \mathcal{W}_2^2(\Omega, \Sigma_i) = c^2 \Sigma^*
\end{equation}
Thus, the optimal barycenter of the scaled task covariances is exactly $c^2 \Sigma^*$.
The scaled Task Arithmetic update is $T'_{\text{TA}} = \sum \lambda_i (c T_i) = c T_{\text{TA}}$, and its covariance is $\Sigma'_{\text{TA}} = c^2 \Sigma_{\text{TA}}$.
The aligned WSA update under the scaled inputs is:
\begin{equation}
  T'_{\text{WSA}} = T'_{\text{TA}} (\Sigma'_{\text{TA}})^{-1/2} (\Sigma'^*)^{1/2}
\end{equation}
Using the properties of matrix powers, for any $M \succ 0$ and $p \in \mathbb{R}$, $(c^2 M)^p = c^{2p} M^p$. Applying this to the inverse square root and square root:
\begin{equation}
  T'_{\text{WSA}} = (c T_{\text{TA}}) (c^2 \Sigma_{\text{TA}})^{-1/2} (c^2 \Sigma^*)^{1/2} = c T_{\text{TA}} \left( c^{-1} \Sigma_{\text{TA}}^{-1/2} \right) \left( c (\Sigma^*)^{1/2} \right) = c T_{\text{TA}} \Sigma_{\text{TA}}^{-1/2} (\Sigma^*)^{1/2} = c T_{\text{WSA}}
\end{equation}
This completes the proof.
\end{proof}

\subsection{Proof of Theorem 4.5 (Bound on Parameter Misalignment and Loss Deviances)}
\label{appendix:proof-loss-bound}
\begin{proof}
Since $L_i(W)$ is $H$-smooth, for any $W, W'$ we have the standard descent lemma:
\begin{equation}
  L_i(W') \leq L_i(W) + \langle \nabla L_i(W), W' - W \rangle + \frac{H}{2} \| W' - W \|_F^2
\end{equation}
Setting $W' = W_{\text{WSA}}$ and $W = W_i$:
\begin{equation}
  L_i(W_{\text{WSA}}) \leq L_i(W_i) + \langle \nabla L_i(W_i), W_{\text{WSA}} - W_i \rangle + \frac{H}{2} \| W_{\text{WSA}} - W_i \|_F^2
\end{equation}
Since $W_i$ is a local minimum, $\nabla L_i(W_i) = 0$. Substituting $W_{\text{WSA}} = W_0 + T_{\text{WSA}}$ and $W_i = W_0 + T_i$ gives:
\begin{equation}
  L_i(W_{\text{WSA}}) - L_i(W_i) \leq \frac{H}{2} \| T_{\text{WSA}} - T_i \|_F^2
\end{equation}
Taking the weighted sum over tasks:
\begin{equation}
  \sum_{i=1}^K \lambda_i \left( L_i(W_{\text{WSA}}) - L_i(W_i) \right) \leq \frac{H}{2} \sum_{i=1}^K \lambda_i \| T_{\text{WSA}} - T_i \|_F^2
\end{equation}
which is exactly $\frac{H}{2} \mathbb{E}_{\lambda} \left[ \| T_{\text{WSA}} - T_i \|_F^2 \right]$.
To analyze the expected Frobenius norm term, we expand the squared Frobenius norm:
\begin{equation}
  \| T_{\text{WSA}} - T_i \|_F^2 = \| T_{\text{WSA}} \|_F^2 + \| T_i \|_F^2 - 2 \text{Tr}\left( T_{\text{WSA}}^\top T_i \right)
\end{equation}
Taking the expectation with respect to $\lambda$:
\begin{equation}
  \mathbb{E}_{\lambda} \left[ \| T_{\text{WSA}} - T_i \|_F^2 \right] = \| T_{\text{WSA}} \|_F^2 + \mathbb{E}_{\lambda} \left[ \| T_i \|_F^2 \right] - 2 \text{Tr}\left( T_{\text{WSA}}^\top \mathbb{E}_{\lambda}[T_i] \right)
\end{equation}
By definition, the expected task update is the Task Arithmetic update $T_{\text{TA}} = \mathbb{E}_{\lambda}[T_i] = \sum \lambda_i T_i$.
For $d_2 \leq d_1$, the row covariance of $T_{\text{WSA}}$ is $\Sigma_{\text{WSA}} = \frac{1}{d_1} T_{\text{WSA}}^\top T_{\text{WSA}} = \Sigma^*$. Thus:
\begin{equation}
  \| T_{\text{WSA}} \|_F^2 = \text{Tr}\left( T_{\text{WSA}}^\top T_{\text{WSA}} \right) = d_1 \text{Tr}(\Sigma_{\text{WSA}}) = d_1 \text{Tr}(\Sigma^*)
\end{equation}
Substituting this back yields:
\begin{equation}
  \mathbb{E}_{\lambda} \left[ \| T_{\text{WSA}} - T_i \|_F^2 \right] = d_1 \text{Tr}(\Sigma^*) + \mathbb{E}_{\lambda} \left[ \| T_i \|_F^2 \right] - 2 \text{Tr}\left( T_{\text{WSA}}^\top T_{\text{TA}} \right)
\end{equation}
This completes the proof.
\end{proof}

\subsection{Proof of Theorem 4.6 (Computational and Memory Complexity Analysis)}
\label{appendix:proof-complexity}
\begin{proof}
We analyze the computational operations and auxiliary memory required for a single layer with base weight matrix $W_0 \in \mathbb{R}^{d_1 \times d_2}$ and $K$ task-specific models.
Let $d_{\min} = \min(d_1, d_2)$ and $d_{\max} = \max(d_1, d_2)$.

\textbf{1. Computational Complexity Derivation:}
The Dual Mode WSA algorithm proceeds through the following computational steps:
\begin{enumerate}
  \item \textbf{Task Update Calculation:} Computing the task-specific updates $T_i = W_i - W_0$ for all $i = 1, \dots, K$ requires element-wise subtraction on $d_1 \times d_2$ matrices. This requires $O(K d_1 d_2)$ operations.
  \item \textbf{Task Covariance Estimation:} Depending on the dimensions, we compute either the row or column covariances:
  \begin{itemize}
    \item If $d_2 \leq d_1$, we compute row covariances $\Sigma_i = \frac{1}{d_1} T_i^\top T_i \in \mathbb{R}^{d_2 \times d_2}$. Each matrix multiplication takes $O(d_1 d_2^2)$ operations, so computing all $K$ covariances takes $O(K d_1 d_2^2) = O(K d_{\max} d_{\min}^2)$ operations.
    \item If $d_1 < d_2$, we compute column covariances $\Sigma_i = \frac{1}{d_2} T_i T_i^\top \in \mathbb{R}^{d_1 \times d_1}$. Storing and computing these takes $O(K d_1^2 d_2) = O(K d_{\max} d_{\min}^2)$ operations.
  \end{itemize}
  Thus, task covariance estimation requires $O(K d_1 d_2 d_{\min})$ operations.
  \item \textbf{Fixed-Point Barycenter Iteration:} For each of the $T$ iterations of the stabilized Alvarez-Esteban algorithm on $d_{\min} \times d_{\min}$ matrices:
  \begin{itemize}
    \item We compute the matrix square root and inverse square root of the current estimate $\Sigma$, which takes $O(d_{\min}^3)$ operations via symmetric eigendecomposition.
    \item For each of the $K$ tasks, we compute the inner product $A = \Sigma^{1/2} \Sigma_i \Sigma^{1/2}$ and its square root $A^{1/2}$. This takes $O(d_{\min}^3)$ operations per task.
    \item The final multiplication and addition of the updated covariance $\Sigma^{(t+1)} = \Sigma_t^{-1/2} (\sum \lambda_i A^{1/2})^2 \Sigma_t^{-1/2}$ takes $O(d_{\min}^3)$ operations.
  \end{itemize}
  Thus, each iteration takes $O(K d_{\min}^3)$ operations, resulting in a total of $O(T K d_{\min}^3)$ operations for $T$ iterations.
  \item \textbf{WSA Alignment Projection:} Computing the Task Arithmetic covariance takes $O(d_{\max} d_{\min}^2)$ operations. Its inverse square root and the barycenter's square root take $O(d_{\min}^3)$ operations. The final alignment projection $T_{\text{WSA}} = T_{\text{TA}} \Sigma_{\text{TA}}^{-1/2} (\Sigma^*)^{1/2}$ (or its column equivalent) involves two matrix multiplications of sizes $(d_{\max} \times d_{\min}) \times (d_{\min} \times d_{\min})$, taking $O(d_{\max} d_{\min}^2)$ operations.
\end{enumerate}
Summing these components, the overall computational complexity per layer is:
\begin{equation}
  O\Big( K d_1 d_2 + K d_1 d_2 d_{\min} + T K d_{\min}^3 + d_1 d_2 d_{\min} \Big) = O\Big( K d_1 d_2 d_{\min} + T K d_{\min}^3 \Big)
\end{equation}
where $d_1 d_2 d_{\min} = d_{\max} d_{\min}^2$.

\textbf{2. Memory Complexity Derivation:}
The auxiliary memory is dominated by storing the task covariance matrices $\{\Sigma_i\}_{i=1}^K$ of size $d_{\min} \times d_{\min}$. This requires storing $K \cdot d_{\min}^2$ floats. Storing the running barycenter $\Sigma$, Task Arithmetic covariance $\Sigma_{\text{TA}}$, and temporary matrix powers takes $O(d_{\min}^2)$ auxiliary memory. The total auxiliary memory complexity is thus $O(K d_{\min}^2)$ per layer, completely independent of the larger dimension $d_{\max}$.

\textbf{3. Comparison with Alternative Frameworks:}
\begin{itemize}
  \item \textbf{Full Parameter Covariance:} Standard full-parameter covariance models (as in classical Laplace approximation merging) treat the entire set of weights as a 1D vector of size $d_1 d_2$. The resulting covariance matrix is of size $(d_1 d_2) \times (d_1 d_2)$. Storing it requires $O(d_1^2 d_2^2)$ memory, and computing its matrix powers requires $O(d_1^3 d_2^3)$ operations. For a modest layer with $d_1 = d_2 = 512$, $d_1 d_2 \approx 2.6 \times 10^5$, making full-parameter models completely intractable. WSA reduces the covariance size to a lightweight $512 \times 512$, taking only 1 MB of memory and milliseconds to solve.
  \item \textbf{OT Fusion:} Discrete optimal transport matching requires solving a network flow or linear programming problem of size $d_{\min} \times d_{\min}$, which takes $O(d_{\min}^3 \log d_{\min})$ operations. Additionally, it requires gathering activation statistics across a substantial evaluation dataset, which incurs significant data collection and forwarding overhead. WSA requires no evaluation data and achieves continuous parameter-space optimal alignment with simple, highly parallelizable matrix operations.
\end{itemize}
\end{proof}

\subsection{Proof of Theorem 4.7 (Robustness and Perturbation Guarantees)}
\label{appendix:proof-robustness}
\begin{proof}
We prove both claims of Theorem 4.7 step-by-step.

\textbf{1. Bures-Wasserstein Metric Lipschitz Bound:}
Recall that the set of zero-mean multivariate Gaussian distributions equipped with the 2-Wasserstein metric is a geodesic space of non-negative curvature, where the metric corresponds exactly to the Bures-Wasserstein distance $d_{\text{BW}}$.
A fundamental property of 2-Wasserstein barycenters in CAT(0) and more general metric spaces is that the barycentric map is a non-expansive mapping (1-Lipschitz) with respect to the product metric.
Specifically, let $\Sigma^*$ be the Bures-Wasserstein barycenter of $\{\Sigma_i\}_{i=1}^K$ under weights $\lambda$, and let $\tilde{\Sigma}^*$ be the Bures-Wasserstein barycenter of the perturbed covariances $\{\tilde{\Sigma}_i\}_{i=1}^K$ under the same weights.
By the non-expansiveness of the barycenter mapping under the 2-Wasserstein metric, we have:
\begin{equation}
  d_{\text{BW}}(\tilde{\Sigma}^*, \Sigma^*) \le \sum_{i=1}^K \lambda_i d_{\text{BW}}(\tilde{\Sigma}_i, \Sigma_i)
\end{equation}
This establishes the first metric Lipschitz bound.

\textbf{2. Frobenius Norm Perturbation Bound:}
We now relate the Bures-Wasserstein metric to the Frobenius norm under the assumption that there exist constants $\alpha, \beta > 0$ such that $\alpha I \le \Sigma_i \le \beta I$ and $\alpha I \le \tilde{\Sigma}_i \le \beta I$ for all $i$.
Since the barycenters $\Sigma^*$ and $\tilde{\Sigma}^*$ lie within the geodesically convex hull of the input matrices in Bures-Wasserstein space, they also satisfy the upper bound $\Sigma^* \le \beta I$ and $\tilde{\Sigma}^* \le \beta I$.
Recall that the Bures-Wasserstein distance between two positive-definite matrices $A$ and $B$ can be written as:
\begin{equation}
  d_{\text{BW}}^2(A, B) = \min_{R \in O(d)} \| A^{1/2} R - B^{1/2} \|_F^2
\end{equation}
Let $R^* \in O(d)$ be the optimal orthogonal alignment matrix achieving this minimum, such that:
\begin{equation}
  d_{\text{BW}}(A, B) = \| A^{1/2} R^* - B^{1/2} \|_F
\end{equation}
We can write the difference $A - B$ algebraically as follows:
\begin{equation}
  A - B = (A^{1/2} R^* - B^{1/2}) (A^{1/2} R^* )^\top + B^{1/2} (A^{1/2} R^* - B^{1/2})^\top
\end{equation}
To verify this identity, we expand the right-hand side:
\begin{align}
  (A^{1/2} R^* - B^{1/2}) R^{*\top} A^{1/2} + B^{1/2} (R^{*\top} A^{1/2} - B^{1/2}) &= A^{1/2} R^* R^{*\top} A^{1/2} - B^{1/2} R^{*\top} A^{1/2} + B^{1/2} R^{*\top} A^{1/2} - B \\
  &= A - B
\end{align}
where we have used the fact that $R^* R^{*\top} = I$ for any orthogonal matrix $R^* \in O(d)$.
Taking the Frobenius norm of both sides of the identity and applying the triangle inequality:
\begin{equation}
  \| A - B \|_F \le \| (A^{1/2} R^* - B^{1/2}) (A^{1/2} R^*)^\top \|_F + \| B^{1/2} (A^{1/2} R^* - B^{1/2})^\top \|_F
\end{equation}
Using the sub-multiplicativity of the Frobenius norm with the spectral norm, we obtain:
\begin{align}
  \| (A^{1/2} R^* - B^{1/2}) (A^{1/2} R^*)^\top \|_F &\le \| A^{1/2} R^* - B^{1/2} \|_F \| A^{1/2} R^* \|_2 \\
  &= d_{\text{BW}}(A, B) \| A^{1/2} \|_2 \\
  &= d_{\text{BW}}(A, B) \sqrt{\|A\|_2}
\end{align}
and similarly:
\begin{align}
  \| B^{1/2} (A^{1/2} R^* - B^{1/2})^\top \|_F &\le \| B^{1/2} \|_2 \| (A^{1/2} R^* - B^{1/2})^\top \|_F \\
  &= \sqrt{\|B\|_2} d_{\text{BW}}(A, B)
\end{align}
Combining these inequalities yields:
\begin{equation}
  \| A - B \|_F \le \left( \sqrt{\|A\|_2} + \sqrt{\|B\|_2} \right) d_{\text{BW}}(A, B)
\end{equation}
Setting $A = \tilde{\Sigma}^*$ and $B = \Sigma^*$, and using the fact that both barycenters satisfy the spectral upper bound $\|\tilde{\Sigma}^*\|_2 \le \beta$ and $\|\Sigma^*\|_2 \le \beta$, we obtain:
\begin{equation}
  \| \tilde{\Sigma}^* - \Sigma^* \|_F \le 2 \sqrt{\beta} d_{\text{BW}}(\tilde{\Sigma}^*, \Sigma*)
\end{equation}
Substituting our first metric Lipschitz bound into this inequality, we have:
\begin{equation}
  \| \tilde{\Sigma}^* - \Sigma^* \|_F \le 2 \sqrt{\beta} \sum_{i=1}^K \lambda_i d_{\text{BW}}(\tilde{\Sigma}_i, \Sigma_i)
\end{equation}
This completes the proof.
\end{proof}

\subsection{Proof of Theorem 4.8 (Low-Rank Preservation)}
\label{appendix:proof-low-rank}
\begin{proof}
Let $\Sigma_1, \dots, \Sigma_K \in \mathcal{P}_d$ be task covariances such that $\text{rank}(\Sigma_i) \le r < d$ for all $i$.
Let $V_i = \text{Range}(\Sigma_i) \subset \mathbb{R}^d$ be the column space (or support) of the matrix $\Sigma_i$.
Let $V = \sum_{i=1}^K V_i$ be the sum of these task subspaces.
The dimension of $V$ is bounded by:
\begin{equation}
  \dim(V) \le \min \left( \sum_{i=1}^K \text{rank}(\Sigma_i), d \right)
\end{equation}
Let $P_V \in \mathbb{R}^{d \times d}$ be the orthogonal projection matrix onto the subspace $V$.
By definition, for each task $i$, the support of $\Sigma_i$ is contained in $V$, which implies that:
\begin{equation}
  P_V \Sigma_i P_V = \Sigma_i, \quad (I - P_V) \Sigma_i = 0
\end{equation}
Now let $\Sigma^*$ be any candidate positive semidefinite matrix for the barycenter.
We can decompose $\Sigma^*$ orthogonally as:
\begin{equation}
  \Sigma^* = \Sigma^*_V + \Sigma^*_{\perp}
\end{equation}
where:
\begin{align}
  \Sigma^*_V &= P_V \Sigma^* P_V \\
  \Sigma^*_{\perp} &= (I - P_V) \Sigma^* (I - P_V)
\end{align}
We analyze the Bures-Wasserstein distance between $\Sigma^*$ and $\Sigma_i$:
\begin{equation}
  d_{\text{BW}}^2(\Sigma^*, \Sigma_i) = \text{Tr}(\Sigma^*) + \text{Tr}(\Sigma_i) - 2 \text{Tr}\left( \left( (\Sigma^*)^{1/2} \Sigma_i (\Sigma^*)^{1/2} \right)^{1/2} \right)
\end{equation}
By the trace identity, $\text{Tr}(\Sigma^*) = \text{Tr}(\Sigma^*_V) + \text{Tr}(\Sigma^*_{\perp})$.
Furthermore, the non-zero eigenvalues of the matrix product $(\Sigma^*)^{1/2} \Sigma_i (\Sigma^*)^{1/2}$ are identical to the non-zero eigenvalues of $\Sigma^* \Sigma_i$. Since $\Sigma_i = P_V \Sigma_i P_V$, we can write:
\begin{equation}
  \Sigma^* \Sigma_i = \Sigma^* P_V \Sigma_i P_V
\end{equation}
By the cyclic permutation property of eigenvalues of matrix products, the non-zero eigenvalues of $(\Sigma^* P_V) (\Sigma_i P_V)$ are identical to those of:
\begin{equation}
  (\Sigma_i P_V) (\Sigma^* P_V) = \Sigma_i (P_V \Sigma^* P_V) = \Sigma_i \Sigma^*_V
\end{equation}
which has the same non-zero eigenvalues as $(\Sigma^*_V)^{1/2} \Sigma_i (\Sigma^*_V)^{1/2}$.
Therefore, the trace of the square root matrix product satisfies:
\begin{equation}
  \text{Tr}\left( \left( (\Sigma^*)^{1/2} \Sigma_i (\Sigma^*)^{1/2} \right)^{1/2} \right) = \text{Tr}\left( \left( (\Sigma^*_V)^{1/2} \Sigma_i (\Sigma^*_V)^{1/2} \right)^{1/2} \right)
\end{equation}
Substituting these back, we obtain:
\begin{align}
  d_{\text{BW}}^2(\Sigma^*, \Sigma_i) &= \text{Tr}(\Sigma^*_V) + \text{Tr}(\Sigma^*_{\perp}) + \text{Tr}(\Sigma_i) - 2 \text{Tr}\left( \left( (\Sigma^*_V)^{1/2} \Sigma_i (\Sigma^*_V)^{1/2} \right)^{1/2} \right) \\
  &= d_{\text{BW}}^2(\Sigma^*_V, \Sigma_i) + \text{Tr}(\Sigma^*_{\perp})
\end{align}
Summing this over all tasks with weights $\lambda$, the Bures-Wasserstein barycenter objective satisfies:
\begin{equation}
  \sum_{i=1}^K \lambda_i d_{\text{BW}}^2(\Sigma^*, \Sigma_i) = \sum_{i=1}^K \lambda_i d_{\text{BW}}^2(\Sigma^*_V, \Sigma_i) + \text{Tr}(\Sigma^*_{\perp})
\end{equation}
Since $\Sigma^*_{\perp}$ is positive semidefinite, $\text{Tr}(\Sigma^*_{\perp}) \ge 0$, with equality if and only if $\Sigma^*_{\perp} = 0$.
Thus, any non-zero out-of-subspace component $\Sigma^*_{\perp}$ strictly increases the objective.
The unique minimizer $\Sigma^*$ of the barycenter objective must therefore satisfy:
\begin{equation}
  \Sigma^*_{\perp} = 0 \implies \Sigma^* = P_V \Sigma^* P_V
\end{equation}
This proves that the column space of the optimal barycenter $\Sigma^*$ is strictly contained in the subspace $V$.
Consequently, the rank of the barycenter is bounded by:
\begin{equation}
  \text{rank}(\Sigma^*) \le \dim(V) \le \min \left( \sum_{i=1}^K \text{rank}(\Sigma_i), d \right)
\end{equation}
This completes the proof.
\end{proof}

\end{document}
